% !TeX program = pdflatex
% papers/naturais.tex — O DEGRAU F_1 → ℕ: a torre binária F_{2w} = F_w + F_w*, e o transporte.
% Convenção do documento: F_w := GF(2^w) — o índice conta BITS, nunca elementos.
%
% Auto-contido: cd papers && pdflatex naturais.tex
\documentclass[11pt,a4paper]{article}
\usepackage[utf8]{inputenc}
\usepackage[T1]{fontenc}
\usepackage[portuguese]{babel}
\usepackage{amsmath,amssymb,amsthm}
\usepackage[margin=2.4cm]{geometry}
\usepackage{booktabs}
\usepackage{xcolor}
\usepackage[hidelinks]{hyperref}

\newtheorem{teorema}{Teorema}
\newtheorem{proposicao}[teorema]{Proposição}
\newtheorem{corolario}[teorema]{Corolário}
\newtheorem{lema}[teorema]{Lema}
\newtheorem{definicao}{Definição}
\newtheorem{observacao}{Observação}

\newcommand{\code}[1]{\texttt{\small #1}}
\newcommand{\dg}{^{\dagger}}
\newcommand{\Word}{\mathcal{W}}
\newcommand{\FF}{\mathbb{F}}
\newcommand{\sepcol}{\quad\penalty0$\big|$\quad\penalty0}
\newcommand{\colunas}[3]{%
  \par\smallskip\noindent{\footnotesize\raggedright
  \textbf{prova} (Algébrico): #1\sepcol
  \textbf{realiza} (aqui): #2\sepcol
  \textbf{verifica}: #3\par}\smallskip}

\definecolor{cinza}{gray}{0.35}
\newcommand{\medido}[1]{\par\medskip\noindent{\small\color{cinza}\textbf{Medido.} #1}\par\medskip}

% ─────────────────────────────────────────────────────────────────────────────
%  papers/gkcapa.tex — a capa do Reino, mínima, para os papers em `article`.
%
%  O estilo.tex completo é do `report` (títulos de capítulo, skak, …). Aqui fica
%  só o que a capa precisa, para o paper continuar a compilar sozinho com
%  pdflatex. O tradutor próprio (tests/tex) carrega o estilo.tex da raiz / o
%  symlink papers/estilo.tex e avalia o mesmo `\gkcapa`.
% ─────────────────────────────────────────────────────────────────────────────
\definecolor{ouro}{HTML}{B8912F}
\definecolor{ouroclaro}{HTML}{F3E9CE}
\definecolor{tinta}{HTML}{1E1B16}
\definecolor{regua}{HTML}{6E6858}
\providecommand{\gknota}{\fontsize{7.62}{11.05}\selectfont}
\providecommand{\gktexto}{\fontsize{10.50}{15.22}\selectfont}
\providecommand{\gksub}{\fontsize{12.33}{17.87}\selectfont}
\providecommand{\gksec}{\fontsize{14.47}{20.98}\selectfont}
\providecommand{\gkcap}{\fontsize{16.99}{24.63}\selectfont}
\providecommand{\gktit}{\fontsize{23.42}{33.95}\selectfont}
\providecommand{\Prj}{\mathbb{P}}
\providecommand{\Trf}{\mathcal{T}}
\providecommand{\gkcapa}[3]{%
\title{\vspace{-0.6cm}
{\color{ouro}\rule{\textwidth}{1.4pt}}\\[6mm]
\textbf{\gktit\color{tinta} Reino Dourado}\\[3mm]{\gksec\itshape\color{regua} Golden Kingdom}\\[8mm]
{\gkcap\scshape\color{ouro} #1}\\[5mm]
{\gksub\color{regua} #2}\\[2mm]
{\gktexto\color{regua}\itshape #3}\\[7mm]
{\color{ouro}\rule{\textwidth}{0.6pt}}\\[9mm]{\gknota\color{regua}\begin{minipage}{\textwidth}\setlength{\parindent}{0pt}%
\textbf{Aviso de Isenção Algorítmica:} O universo, as leis e as entidades contidas nestas páginas ---
incluindo felinos matriciais, aves de rapina algébricas e amuletos de controle de fluxo --- são
construções de pura ficção formal. Esta obra não descreve a realidade física, não serve como
engenharia reversa do cosmos e não constitui doutrina religiosa. Qualquer semelhança com bugs reais do seu
sistema operacional é mera coincidência (ou falta de reza). Prossiga por sua própria conta e
risco.\end{minipage}}}
\author{}\date{}}


\begin{document}
\gkcapa{Os naturais a partir do binário}
       {$F_{2w}=F_w+F_w^{*}$ --- quatro dobras de um bit ao byte,
        e $\mathbb{N}$ é a torre mais o transporte}
       {oito leis, dobra dual, encaixe por prefixos, base ortonormal, transporte}
\maketitle

%======================================================================
\begin{abstract}
\textbf{Convenção, antes de tudo:} $\FF_2$ é o corpo de dois elementos e
\[
\boxed{\;F_w:=\mathrm{GF}(2^{w})\;}\qquad\text{logo}\qquad
F_1=\FF_2,\quad \lvert F_w\rvert=2^{w},\quad \lvert F_8\rvert=256 .
\]
O índice conta \textbf{bits}, nunca elementos --- e é o mesmo índice em todo o
documento.

\medskip \textbf{Primeiro degrau da construção.} Começa-se no bit e sobe-se:
o que aqui se entrega é o suporte sobre o qual o \code{inteiros.tex} irá
duplicar e quocientar ($\mathbb{N}^{2}/{\sim}=\mathbb{Z}$). Aqui duplica-se o
\textbf{corpo} e \textbf{não} se quocienta:
$F_{2w}=F_w\oplus\sigma F_w$ com $\sigma^{2}=\sigma+\lambda$. Partindo de
$F_1=\FF_2$ --- um bit, duas portas --- quatro andares dão
$F_1\to F_2\to F_4\to F_8$, de cardinais $2,4,16,256$; e o suporte do topo é
exactamente $\Word_8=\{0,\ldots,255\}$, o envelope do \code{inteiros.tex}
(§\ref{sec:encaixe}). A régua $\lambda_k$ não se escreve: deriva-se, e o
que a busca devolve é $2^{2^{k}-1}$ --- a régua do \emph{nim} de Conway
(§\ref{sec:conway}). A base dos produtos dos três geradores é
$e_k=2^{k}$ --- a base ortonormal das oito leis (§\ref{sec:base}). E o que separa
o corpo de $\mathbb{N}$ tem nome e conta-se: o \textbf{transporte}
(§\ref{sec:transporte}).
\end{abstract}

%======================================================================
\section*{Onde este documento vive}

A construção lê-se num sentido só --- do bit para $\mathbb{R}$ ---, e este é
o seu \textbf{primeiro} degrau ($k=-1$ na numeração da casa, porque a casa já
contava a partir de $\mathbb{N}\!\to\!\mathbb{Z}$):

\[
\boxed{\;
\underbrace{\FF_2}_{\text{um bit}}
\;\xrightarrow[\text{3 dobras duais}]{\ \code{naturais}\ }\;
\underbrace{F_8\equiv\Word_8}_{\{0,\ldots,255\}}
\;\xrightarrow[\text{TRANSPORTE}]{}\;
\mathbb{N}
\;\xrightarrow[\text{soma}]{\ \code{inteiros}\ }\;
\mathbb{Z}
\;\xrightarrow[\text{produto}]{\ \code{racionais}\ }\;
\mathbb{Q}
\;\xrightarrow[\text{corte}]{\ \code{reais}\ }\;
\mathbb{R}\;}
\]
\noindent{\footnotesize Lê-se da esquerda para a direita: é o sentido em que se
\textbf{constrói}. Cada paper recebe o objecto do degrau à sua esquerda e
entrega o do degrau à sua direita; nenhum redefine o que recebeu.\par}

\noindent A cadeia completa, comum aos três papers
--- a mesma nos três:

\begin{center}\small
\begin{tabular}{@{}clllll@{}}
\toprule
$k$ & degrau & ganha volta & paga & dual do andar (e o seu Fix) & paper \\
\midrule
$-1$ & $F_1\!\to\!F_2\!\to\!F_4\!\to\!F_8$ & a \textbf{largura} (3 dobras)
     & a \textbf{ordem} (car.\ $2$) & $x^{*}$ (2.ª raiz); Fix $=$ andar de baixo
     & \code{naturais.tex} \\
$0$  & $\mathbb{N}\!\to\!\mathbb{Z}$ & a \textbf{soma}
     & o sinal e a boa ordem & $\iota(a,b)=(b,a)\Rightarrow$ oposto; Fix $=\{0\}$
     & \code{inteiros.tex} \\
$1$  & $\mathbb{Z}\!\to\!\mathbb{Q}$ & a \textbf{multiplicação}
     & $0^{-1}$ não existe & $\iota(a,b)=(b,a)\Rightarrow$ inverso; Fix $=\{\pm1\}$
     & \code{racionais.tex} \\
$2$  & $\mathbb{Q}\!\to\!\mathbb{R}$ & o \textbf{corte} (completação)
     & a \textbf{enumerabilidade} & $(A,B)\mapsto(B,A)$; Fix: o corte que fecha
     & \code{reais.tex} \\
\bottomrule
\end{tabular}
\end{center}
\noindent{\footnotesize A numeração é a de \code{arquitetura sec:descida},
onde $k=0$ é $\mathbb{N}\!\to\!\mathbb{Z}$; a torre binária é \emph{anterior} a
esse degrau, e por isso $k=-1$. O passo é sempre $T_{k+1}=T_k+T_k^{*}$
(\code{analitico thm:tecidos}). \textbf{Entre $k=-1$ e $k=0$ não há dobra:
há o TRANSPORTE} --- $a+b=(a\oplus b)+2(a\wedge b)$
(\code{naturais.tex thm:transporte}) ---, e é ele que leva o suporte $\Word_8$
do topo da torre binária às operações de $\mathbb{N}$. Nos degraus $k\ge0$
duplica-se e \textbf{quocienta-se}; em $k=-1$ duplica-se e \textbf{não} se
quocienta (a segunda cópia é multiplicada por $\sigma$).\par}

\noindent E o dual, que em cada degrau tem \emph{duas} partes e não se
confundem:

\[
\boxed{
\begin{array}{lll}
 & \textbf{dual do CENTRO (Lei 1)} & \textbf{dual do ANDAR (o que o degrau acrescenta)}\\[3pt]
k=-1 & \nu(x)=-x=x\ \text{(degenera)} & x^{*}=(a_0{+}a_1)+a_1\sigma\ \Rightarrow\ T=x\oplus x^{*},\ N=x\otimes x^{*}\\[2pt]
k=0  & \nu(z)=-z & \iota(a,b)=(b,a)\ \Rightarrow\ \text{oposto aditivo}\\[2pt]
k=1  & \nu(a,b)=(-a,b) & \iota(a,b)=(b,a)\ \Rightarrow\ \text{inverso multiplicativo}
\end{array}}
\]
\noindent{\footnotesize \textbf{O mesmo \emph{swap} induz o oposto num degrau e
o inverso no seguinte} --- é o dado que muda de papel, não a fórmula. E em
$k=-1$ o dual do centro degenera ($-x=x$), pelo que o dual tem de vir de outro
sítio: da segunda raiz de $\sigma^2+\sigma+\lambda$. Nenhuma destas involuções
é o dual linear $V\!\cong\!V^{**}$ da Lei~2, nem o dual quadrático
$s\dg=-1/s$ do \emph{Corpo Algébrico}: são quatro noções, e não se confundem.\par}
\textbf{Molde dual} (Lei~4, $T+T^{*}$): duplicar, ler nos dois lados. A
diferença face aos degraus de cima é que aqui \textbf{não há quociente}: a
segunda cópia não é identificada, é \emph{multiplicada} por $\sigma$. O que se
paga por andar é a régua $\lambda$; o que se ganha é a largura, que dobra.

\medskip\noindent\textbf{O que este degrau ENTREGA.} Ao degrau seguinte
(\code{inteiros.tex}, $k=0$) entrega três coisas: o \textbf{suporte}
$\Word_8=\{0,\ldots,255\}$ com o encaixe por prefixos (Thm.~\ref{thm:encaixe});
a \textbf{soma} de $\mathbb{N}$ reconstruída a partir das duas primitivas de um
bit (Thm.~\ref{thm:transporte}); e a \textbf{base ortonormal} de oito que
autoriza a operação bit a bit (Thm.~\ref{thm:base}). O que \emph{não} entrega,
e o degrau seguinte terá de comprar, é a ordem.

\medskip\noindent\textbf{A leitura em uma frase.} A subtracção, que o
\code{inteiros.tex} compra ao preço da ordem, aqui já é total no chão
($-x=x$); e a \emph{ordem} é que fica por comprar. É o mesmo par visto do outro
lado: \emph{a face que opera não ordena; a que ordena não opera}
(\code{algebrico}, §O byte).

\medskip\noindent\textbf{O que este documento RECEBE e o que ACRESCENTA:}

\[
\boxed{
\begin{array}{lll}
\textbf{RECEBE do centro} & \textbf{onde está lá} & \textbf{o que se faz aqui}\\[3pt]
\text{a torre }T_{k+1}=T_k+T_k^{*} & \code{analitico thm:tecidos} & \text{instancia o degrau }F_1\to F_8\\[2pt]
\text{o byte é corpo, posição }k=\text{Lei }k & \code{algebrico §O byte} & \text{constrói-o por dobras, não por }f_8\\[2pt]
\text{a base das oito leis é ortonormal} & \code{ortonormal.c §O0} & \text{obtém-na como PRODUTOS dos geradores}\\[2pt]
\text{o construtor }(K,m)\mapsto K[\sigma_m] & \code{algebrico thm:autossimilar} & \text{o mesmo, em característica }2\\[2pt]
\text{o envelope }\Word_8 & \code{inteiros def:w8} & \text{mostra que }\Word_8\text{ é o SUPORTE de }F_8\\[6pt]
\textbf{ACRESCENTA} & \multicolumn{2}{l}{\text{---e é só isto---}}\\[3pt]
\text{o ENCAIXE por prefixos} & \text{Thm.~\ref{thm:encaixe}} & \{0,\ldots,2^{2^k}-1\}\text{ é subcorpo}\\[2pt]
\text{a RÉGUA derivada} & \text{Prop.~\ref{prop:regua}} & \lambda_k=2^{2^{k}-1}\text{, e é a do \emph{nim}}\\[2pt]
\text{a BASE da torre} & \text{Thm.~\ref{thm:base}} & e_k=2^k\text{ pelos três geradores}\\[2pt]
\text{o TRANSPORTE} & \text{Thm.~\ref{thm:transporte}} & a+b=(a\oplus b)+2(a\wedge b)\\[2pt]
\text{o ESCOPO} & \text{§\ref{sec:escopo}} & F_8\text{ é o suporte de }\Word_8,\ \text{não }\mathbb{N}
\end{array}}
\]

%======================================================================
\section{As oito leis --- marcos da construção}\label{sec:leis}

\noindent Onde cada lei vive na cadeia --- tabela comum aos três papers
--- a mesma nos três, para que os estatutos não divirjam entre eles:

\begin{center}\footnotesize
\begin{tabular}{@{}cp{2.75cm}p{2.75cm}p{2.75cm}p{2.75cm}@{}}
\toprule
lei & $k{=}-1$: $F_1\!\to\!F_8$ & $k{=}0$: $\mathbb{N}\!\to\!\mathbb{Z}$ & $k{=}1$: $\mathbb{Z}\!\to\!\mathbb{Q}$ & $k{=}2$: $\mathbb{Q}\!\to\!\mathbb{R}$ \\
\midrule
\textbf{0} divisão do zero
 & $x\oplus x=0$; só o $0$ sem inversa \hfill\textsc{expl.}
 & subtracção parcial; folha $[(1,0)]\oplus[(0,1)]=0$ \hfill\textsc{expl.}
 & $0^{-1}$ não existe; $0\!\leftrightarrow\!\infty$ \hfill\textsc{mista}
 & o \textbf{buraco}: corte sem racional em cima \hfill\textsc{expl.} \\
\textbf{1} dual ($1\dg=-1$)
 & \emph{degenera}: $-x=x$ \hfill\textsc{expl.}
 & $\iota$ no par; $\nu(z)=-z$ \hfill\textsc{expl.}
 & $\nu(a,b)=(-a,b)$ \hfill\textsc{expl.}
 & $\nu(x)=-x$; do andar: $(A,B)\!\mapsto\!(B,A)$ \hfill\textsc{expl.} \\
\textbf{2} bidual
 & $x^{**}=x$; automorfismo \hfill\textsc{expl.}
 & Dir/Cruz; $R^{4}=\mathrm{id}$ \hfill\textsc{mista}
 & $V\cong V^{**}$ canónico \hfill\textsc{canón.}
 & $B=\mathbb{Q}\setminus A$: complementar $2\times$ \hfill\textsc{expl.} \\
\textbf{3} trial $\{-1,0,+1\}$
 & \emph{degenera}: sem cúspide \hfill\textsc{expl.}
 & cúspide $\tau=0$; $+0=-0$ \hfill\textsc{expl.}
 & $\operatorname{sgn}(a/b)$ \hfill\textsc{expl.}
 & $\tau=0$ é \emph{vazio} $\iff$ irracional \hfill\textsc{expl.} \\
\textbf{4} tetral $T+T^{*}$
 & $F_{2w}=F_w\oplus\sigma F_w$ \hfill\textsc{expl.}
 & $\mathbb{N}^{2}/{\sim}$; classe $[(a,b)]$ \hfill\textsc{expl.}
 & recebida (\code{thm:juntas}) \hfill\textsc{receb.}
 & $(A,B)$, e quocientar aberto/fechado \hfill\textsc{expl.} \\
\textbf{5} pental $x^{2}=-1$
 & \emph{degenera}: uma raiz só \hfill\textsc{expl.}
 & $G_{\mathrm{geom}}\bmod2$ \hfill\textsc{receb.}
 & $J(a,b)=(-b,a)$, $J^{2}=-I$ \hfill\textsc{expl.}
 & \emph{não existe}: ordenado \hfill\textsc{expl.} \\
\textbf{6} hexal (interface)
 & base ortonormal: coordenada $=$ ler o bit \hfill\textsc{receb.}
 & $\otimes\leftrightarrow(a-b)(c-d)$ \hfill\textsc{receb.}
 & interface \hfill\textsc{receb.}
 & densidade de $\mathbb{Q}$ em $\mathbb{R}$ \hfill\textsc{receb.} \\
\textbf{7} octonião dual
 & encaixe por prefixos \hfill\textsc{expl.}
 & $q\parallel p$; par $\neq$ classe; $\Word_8$ \hfill\textsc{expl.}
 & junta dual \hfill\textsc{receb.}
 & três realizações ligadas sem fundidas \hfill\textsc{expl.} \\
\bottomrule
\end{tabular}
\end{center}
\noindent{\footnotesize \textsc{expl.} realizada neste degrau ·
\textsc{receb.} interface herdada do centro, não derivada aqui ·
\textsc{mista}/\textsc{canón.} ver o paper do degrau.
\textbf{Três leis degeneram em $k=-1$} (1, 3, 5) --- e é isso que permite a
torre binária dobrar sem excepção. A Lei~5 falta nos dois extremos por razões
\emph{opostas}: em $k=-1$ por característica~2, em $k=2$ por ordem.
\textbf{Nenhum degrau realiza as oito}; a coluna diz onde cada uma vive.\par}

\noindent E, neste degrau, com o § onde cada uma se realiza:

\begin{center}\small
\begin{tabular}{@{}clp{5.0cm}cl@{}}
\toprule
 & lei (centro) & realização $F_1\!\to\!F_8$ & estatuto & § \\
\midrule
\textbf{0} & divisão do zero & $x\oplus x=0$ imediato; $0$ sem inversa; $\mathbb{P}^1(\FF_2)$ com $0\dg=\infty$ & explícita & \ref{sec:dobra} \\
\textbf{1} & dual ($1\dg=-1$) & degenera: $-x=x$; o dual estrutural é $x^{*}$ (a 2.ª raiz) & explícita & \ref{sec:dual} \\
\textbf{2} & bidual ($T\dg=-T$) & $x^{**}=x$; $x^{*}$ automorfismo nos dois lados & explícita & \ref{sec:dual} \\
\textbf{3} & trial $\{-1,0,+1\}$ & colapsa em $\{-1,+1\}$: a cúspide não existe em char.~$2$ & explícita & \ref{sec:trial} \\
\textbf{4} & tetral ($T+T^{*}$) & $F_{2w}=F_w\oplus\sigma F_w$, $\sigma^2=\sigma+\lambda$ & explícita & \ref{sec:dobra} \\
\textbf{5} & pental ($x^{2}=-1$) & degenera: $x^2=1\Rightarrow x=1$; o Frobenius é bijecção & explícita & \ref{sec:dual} \\
\textbf{6} & hexal (interface) & a base ortonormal: coordenada $=$ produto interno $=$ ler o bit & recebida & \ref{sec:base} \\
\textbf{7} & octonião dual & encaixe: os andares ligam-se sem se fundirem & explícita & \ref{sec:encaixe} \\
\bottomrule
\end{tabular}
\end{center}

\noindent Três leis \textbf{degeneram} neste degrau (1, 3 e 5) --- e a degenerescência
não é perda de conteúdo: é o que \emph{permite} a torre dobrar sem excepção
(Prop.~\ref{prop:trial}). O \code{inteiros.tex} tinha duas leis
\emph{recebidas} por não se derivarem no seu degrau; aqui só a Lei~6 o é.

%======================================================================
\section{A dobra --- Leis 0 e 4}\label{sec:dobra}

\begin{definicao}[o construtor]\label{def:dobra}
Seja $F_w=\mathrm{GF}(2^{w})$ e $\lambda\in F_w$. O andar seguinte é
\[
\boxed{\;F_{2w}=F_w\oplus\sigma F_w,\qquad \sigma^{2}=\sigma+\lambda\;}
\]
--- o \emph{mesmo} molde $T+T^{*}$ do \code{inteiros.tex}, sem quociente. Escrito
nas primitivas, com $x=a_0+a_1\sigma$ e $y=b_0+b_1\sigma$:
\[
x\oplus y=(a_0\oplus b_0)+(a_1\oplus b_1)\sigma,
\qquad
x\otimes y=(a_0b_0\oplus\lambda a_1b_1)+(a_0b_1\oplus a_1b_0\oplus a_1b_1)\sigma .
\]
A \textbf{representação} é a que já existe: $a_0+a_1\sigma$ guarda-se como o
inteiro $a_0+2^{w}a_1$, onde $w$ é a largura de $F_w$ em bits. Nenhuma codificação se
escolhe --- as coordenadas \emph{são} os bits.
\end{definicao}

\noindent\textbf{Lei~0} neste degrau: em $\FF_2$ a soma é o XOR e o oposto é o
próprio elemento, logo a folha $x\oplus x\dg=0$ do \code{inteiros.tex} é
\emph{imediata} e a subtracção já é total. O que resta indivisível é o zero:
$0$ não tem inversa, e é ele o único ponto onde a fibra falta
(\code{algebrico sec:lei0}); em $\mathbb{P}^1(\FF_2)$ ele troca com $\infty$
(\code{umbit.c} §B1).

\begin{proposicao}[a régua deriva-se, e é a metade]\label{prop:regua}
$\sigma^{2}+\sigma+\lambda$ é irredutível sobre $F_w$ se e só se
$y^{2}+y\neq\lambda$ para todo $y\in F_w$ --- e isso vale para exactamente
metade dos $\lambda$ (os de traço $1$). Procurando o \textbf{menor} $\lambda$
que serve em cada andar obtém-se
\[
\boxed{\;\lambda_0=1,\quad\lambda_1=2,\quad\lambda_2=8,
\qquad\text{isto é}\qquad \lambda_k=2^{2^{k}-1}=\tfrac{1}{2}\lvert F_{2^{k}}\rvert\;}
\]
--- metade da potência de Fermat que abre o andar. \textbf{Estatuto:} que este
$\lambda$ \emph{serve} é demonstrado abaixo e reaparece como consequência das
regras do \emph{nim} (Thm.~\ref{thm:nim}); que é o \textbf{menor} que serve está
\emph{medido} nos três andares varridos (§NB0), não demonstrado para $k\ge3$.
\end{proposicao}

\begin{proof}
Um polinómio de grau $2$ é redutível sobre $F_w$ exactamente quando tem raiz em
$F_w$; $\sigma^2+\sigma+\lambda$ tem raiz $y$ sse $y^2+y=\lambda$. A aplicação
$\wp(y)=y^2+y$ é aditiva (char.~$2$) com núcleo $\{0,1\}$, logo a sua imagem
tem índice $2$: metade dos $\lambda$ está na imagem (redutíveis) e metade fora.
O valor de $\lambda_k$ é o resultado da busca --- \emph{medido}, não escrito
(\code{naturais.c} §NB0, §NB10). $\square$
\end{proof}

\begin{observacao}[nenhuma constante entra por cima]
A régua é o único parâmetro do construtor, e é o único sítio onde uma constante
podia entrar à mão. Ela deriva-se por varrimento exaustivo do andar de baixo;
o valor $2^{2^k-1}$ é \emph{conclusão}, não entrada. Compare-se com a outra
régua possível --- o polinómio $x^8+x^4+x^3+x+1$ do AES ---, que dá corpo
igualmente legítimo e \textbf{não} encaixa (§\ref{sec:escopo}).
\end{observacao}

\colunas{\code{algebrico thm:autossimilar} ($(K,m)\mapsto K[\sigma_m]$)}
        {$\sigma^2=\sigma+\lambda$ em char.\ $2$; $\lambda$ derivado}
        {\code{naturais.c} §NB0 ($4/4$ andares, toda inversa exibida); §NB5}

%======================================================================
\section{O dual --- Leis 1, 2 e 5}\label{sec:dual}

\begin{teorema}[a conjugação é o dual, e é o Frobenius do andar]\label{thm:dual}
As duas raízes de $\sigma^2+\sigma+\lambda$ são $\sigma$ e $\sigma+1$. A troca
\[
\boxed{\;(a_0+a_1\sigma)^{*}=(a_0+a_1)+a_1\sigma\;}
\]
satisfaz: \textbf{(i)} $x^{**}=x$ (Lei~2); \textbf{(ii)} é automorfismo dos
\emph{dois} lados, $(x\oplus y)^{*}=x^{*}\oplus y^{*}$ e
$(x\otimes y)^{*}=x^{*}\otimes y^{*}$; \textbf{(iii)}
$x^{*}=x^{\lvert F\rvert}$ --- é o Frobenius relativo ao andar de baixo;
\textbf{(iv)} $\{x:x^{*}=x\}=F$, isto é, o conjunto de pontos fixos do dual
\emph{é} o andar de baixo. Traço e norma descem:
$T(x)=x\oplus x^{*}=a_1$ e $N(x)=x\otimes x^{*}\in F$.
\end{teorema}

\begin{proof}
(i) e (ii) por cálculo directo nas fórmulas da Def.~\ref{def:dobra}; (iii)
porque o grupo de Galois de uma extensão quadrática de corpos finitos é
gerado pelo Frobenius $x\mapsto x^{\lvert F\rvert}$ e o único automorfismo
não trivial é a troca das raízes; (iv) $x^{*}=x$ força $a_1=0$.
\emph{Verifica:} \code{naturais.c} §NB2 --- involução, automorfismo em soma e
produto, e ponto fixo $=$ prefixo, exaustivo nos três andares, zero falhas.
$\square$
\end{proof}

\noindent\textbf{Lei~1 e Lei~5 degeneram, e diz-se como.} $1\dg=-1$ e $-1=1$:
o dual da Lei~1 é a identidade. $x^{2}=-1$ torna-se $x^2=1$, que em
característica $2$ tem \emph{uma} só solução, porque $x\mapsto x^2$ é bijecção
(o Frobenius). Não há $i$ neste corpo --- e é exactamente por isso que o dual
tem de vir da \emph{segunda raiz de uma equação nova} (Thm.~\ref{thm:dual}) e
não da raiz de $-1$. O dual não desaparece: muda de sítio.

\colunas{\code{analitico sub:oitoleis}; \code{algebrico thm:corpo-dual}}
        {$x^{*}$ (2.ª raiz); $T$, $N$; degenerescência de 1 e 5}
        {\code{naturais.c} §NB1, §NB2; \code{umbit.c} §B0}

%======================================================================
\section{O encaixe --- Lei 7}\label{sec:encaixe}

\begin{teorema}[o suporte é o prefixo de $\mathbb{N}$]\label{thm:encaixe}
Com a representação da Def.~\ref{def:dobra}, o suporte do andar $k$ é
exactamente $\{0,1,\ldots,2^{2^{k}}-1\}$, e para $k<k'$ as operações do andar
$k'$ \textbf{restringem} às do andar $k$:
\[
\boxed{\;\underbrace{\{0,1\}}_{F_1}\subset\underbrace{\{0,\ldots,3\}}_{F_2}
\subset\underbrace{\{0,\ldots,15\}}_{F_4}\subset\underbrace{\{0,\ldots,255\}}_{F_8}=\Word_8\;}
\]
--- cadeia de subcorpos, sem conversão de representação entre andares.
\end{teorema}

\begin{proof}
Indução. A soma é o XOR em todos os andares e não mexe em bits acima de $w$.
Para o produto, sejam $x,y$ no andar $k'$ com $a_1=b_1=0$; então nas fórmulas
da Def.~\ref{def:dobra} o coeficiente de $\sigma$ é
$a_0b_1\oplus a_1b_0\oplus a_1b_1=0$ e a parte baixa é $a_0b_0$, calculada
no andar de baixo. Logo o prefixo é fechado e a operação induzida é a do andar
$k$. Por hipótese de indução, desce até $\FF_2$.
\emph{Verifica:} \code{naturais.c} §NB3 --- $276$ pares, zero falhas; e
§NB4 pelo \emph{outro} caminho: o conjunto $\{x\in F_8:x^{2^{w}}=x\}$, que
não sabe nada da construção, tem $2^{w}$ elementos e é o mesmo prefixo, para
$w=1,2,4,8$. $\square$
\end{proof}

\begin{corolario}[$\Word_8$ é o suporte de $F_8$]\label{cor:w8}
O envelope $\Word_8=\{0,\ldots,255\}$ do \code{inteiros.tex} (Def.~\code{def:w8})
não é uma restrição imposta de fora aos naturais: é o \textbf{suporte} do topo
desta torre, e a identificação é a identidade --- o bit $j$ do inteiro é a
coordenada $j$ na base. A cadeia
$\Word_8\to E_{16}\to\cdots\to\mathrm{GF}(2)$ da \code{arquitetura sec:descida}
percorre esta mesma torre \emph{ao contrário}.
\end{corolario}

\noindent\textbf{Lei~7} literal: ligar sem fundir. O par $(a_0,a_1)$ continua
legível dentro do byte depois de multiplicado --- os andares partilham suporte
sem se confundirem, e nenhum \emph{cast} acontece na fronteira.

%======================================================================
\section{O trial em característica 2 --- Lei 3}\label{sec:trial}

\begin{proposicao}[a cúspide não existe aqui]\label{prop:trial}
Sobre um corpo de característica $2$, $\sigma^2+\sigma+\lambda$ tem $0$ ou $2$
raízes --- \textbf{nunca} $1$. Logo $\tau\in\{-1,+1\}$ e o caso $\tau=0$
(\code{algebrico thm:cuspide}) é vazio neste degrau.
\end{proposicao}

\begin{proof}
Raiz dupla exigiria anular a derivada formal, que é a constante $1$. Ou, sem
derivadas: se $y$ é raiz, $y+1$ também é, pois
$(y+1)^2+(y+1)+\lambda=y^2+1+y+1+\lambda=y^2+y+\lambda$; e $y\neq y+1$ porque
$1\neq0$. Logo as raízes vêm aos pares.
\emph{Verifica:} \code{naturais.c} §NB5 --- varridos todos os $\lambda$ dos
três andares: $11$ com zero raízes, $11$ com duas, \textbf{zero} com uma. $\square$
\end{proof}

\begin{observacao}[por que isto importa]
Nos degraus de cima a cúspide é o caso que exige tratamento à parte --- é onde
as duas folhas colidem. A dobra binária não tem esse caso: cada $\lambda$ ou
serve (duas folhas) ou não serve (nenhuma), e metade serve sempre
(Prop.~\ref{prop:regua}). \emph{É por o trial degenerar que a torre dobra sem
excepção} --- a ausência do caso do meio é o que torna a construção uniforme
nos quatro andares.
\end{observacao}

%======================================================================
\section{A base de oito, e é a ortonormal --- Lei 6}\label{sec:base}

\begin{lema}[multiplicar por $\sigma$ é deslocar]\label{lem:desloc}
Seja $\sigma$ o gerador da dobra $F_w\to F_{2w}$, isto é, o elemento
representado pelo inteiro $2^{w}$. Então, para todo $x\in F_w$ (ou seja,
$x<2^{w}$ como inteiro),
\[
x\otimes\sigma=2^{w}x\qquad\text{(o produto \emph{ordinário} em }\mathbb{N}).
\]
\end{lema}

\begin{proof}
Na Def.~\ref{def:dobra}, $x=(x,0)$ e $\sigma=(0,1)$. A parte baixa do produto é
$a_0b_0\oplus\lambda a_1b_1=x\cdot0\oplus\lambda\cdot0\cdot1=0$ e o coeficiente
de $\sigma$ é $a_0b_1\oplus a_1b_0\oplus a_1b_1=x$. Logo
$x\otimes\sigma=(0,x)$, que se guarda como o inteiro $0+2^{w}x$. $\square$
\end{proof}

\begin{teorema}[a base dos produtos dos três geradores é $e_k=2^{k}$]\label{thm:base}
Sejam $\sigma_0=2$, $\sigma_1=4$, $\sigma_2=16$ os geradores das três dobras
(o primeiro elemento novo de cada andar). Então, para $k=0,\ldots,7$,
\[
\boxed{\;e_k\;=\;\prod_{j:\,k_j=1}\sigma_j\;=\;2^{k},
\qquad k=k_0+2k_1+4k_2\;}
\]
--- \emph{o índice do vector é o subconjunto de dobras que o produziu, lido em
binário}. Esta é a base $\{1,2,4,\ldots,128\}$ que \code{ortonormal.c} §O0
mostra ser \textbf{ortonormal} para $\langle a,b\rangle=$ paridade$(a\wedge b)$,
a única forma bilinear de $\FF_2$: $\langle e_i,e_j\rangle=\delta_{ij}$.
\end{teorema}

\begin{proof}
\textbf{(i) É base.} Pela multiplicatividade do grau numa torre de extensões:
se $B$ é base de $L/K$ e $C$ é base de $M/L$, então $\{bc\}$ é base de $M/K$.
Aplicada três vezes a
$\FF_2\subset F_2\subset F_4\subset F_8$ com as bases relativas
$\{1,\sigma_0\}$, $\{1,\sigma_1\}$, $\{1,\sigma_2\}$, dá base de $F_8$ sobre
$\FF_2$ com $2^{3}=8$ elementos, que são os produtos
$\sigma_0^{k_0}\sigma_1^{k_1}\sigma_2^{k_2}$.
\emph{Não} se trata de um produto tensorial de corpos --- as extensões são
\textbf{encaixadas}, não independentes, e de facto
$F_2\otimes_{\FF_2}F_2$ nem sequer é corpo. O que se usa é só a
multiplicatividade do grau.

\textbf{(ii) O valor é $2^{k}$.} Por indução no número de factores, ordenados
por $j$ crescente. O produto parcial dos factores de índice $<j$ pertence a
$F_{2^{j}}$ --- ou seja, é um inteiro $<2^{2^{j}}$ ---, porque $\sigma_i$ com
$i<j$ vive no andar $F_{2^{j}}$. Multiplicá-lo por $\sigma_j$ é então, pelo
Lema~\ref{lem:desloc}, deslocá-lo $2^{j}$ bits. Logo
\[
e_k=2^{k_0}\cdot2^{2k_1}\cdot2^{4k_2}=2^{\,k_0+2k_1+4k_2}=2^{k}.
\]
Nenhum passo reduz módulo $\sigma_j^2=\sigma_j+\lambda_j$: a redução só entra
quando um gerador se multiplica \emph{por si próprio}, e nos $e_k$ cada
gerador aparece no máximo uma vez.

\textbf{(iii)} Gram $=I$ é \code{ortonormal.c} §O0.
\emph{Verifica:} \code{naturais.c} §NB11 --- os oito produtos, a Gram e a
filtração, zero falhas. $\square$
\end{proof}

\noindent\textbf{Terminologia, para não haver briga.} Chama-se a
$\{1,2,4,\ldots,128\}$ a \emph{base dos produtos dos geradores} --- é uma base
\textbf{vectorial} de $F_8$ sobre $\FF_2$, cujos elementos se obtêm
multiplicativamente. \emph{Não} se invoca aqui a noção técnica de «base
multiplicativa» de um corpo.

\begin{corolario}[a coordenada é ler o bit, e o encaixe é a filtração]
Porque a Gram é a identidade, a coordenada recupera-se por
$a_k=\langle a,e_k\rangle$ --- um AND e uma paridade, isto é, \emph{ler o bit};
numa base qualquer seria preciso inverter a matriz de Gram
(\code{ortonormal.c} §O2, §O4). E os $w$ primeiros vectores da base geram
$F_w$: o encaixe do Thm.~\ref{thm:encaixe} é a \textbf{filtração} desta
base. \emph{O \code{uint8\_t} não é conveniência de máquina: é a base
ortonormal escrita em hardware}, e agora também a base que a torre produz
pelos seus geradores.
\end{corolario}

%======================================================================
\section{$\mathbb{N}$ é a torre mais o transporte}\label{sec:transporte}

\noindent Chegados ao topo, o suporte já é $\Word_8$ --- mas as operações são
as do corpo, não as de $\mathbb{N}$. Falta exactamente uma coisa, e ela tem nome.

\begin{lema}[a decomposição]\label{lem:carry}
Para todos $a,b\in\mathbb{N}$:
\[
\boxed{\;a+b=(a\oplus b)+2\,(a\wedge b)\;}
\]
\end{lema}

\begin{proof}
Posição a posição: para $a_j,b_j\in\{0,1\}$ verifica-se nos quatro casos que
$a_j+b_j=(a_j\oplus b_j)+2(a_j\wedge b_j)$ em $\mathbb{N}$. Multiplicando por
$2^{j}$ e somando em $j$:
$a+b=\sum_j 2^{j}(a_j+b_j)=(a\oplus b)+2(a\wedge b)$. $\square$
\end{proof}

\begin{teorema}[$\mathbb{N}$ do binário]\label{thm:transporte}
A iteração
\[
(a,b)\longmapsto\bigl(a\oplus b,\;2\,(a\wedge b)\bigr)
\]
preserva a soma (Lema~\ref{lem:carry}) e \textbf{converge em tempo finito}: se
$a,b<2^{n}$, ao fim de no máximo $n+1$ passos o par é $(a+b,\,0)$,
\[
\boxed{\;(a,b)\longmapsto(a\oplus b,\;2(a\wedge b))\longmapsto\cdots
\longmapsto(a+b,\;0)\;}
\]
--- o transporte a ser propagado, literalmente. Os pontos fixos da iteração são
exactamente os pares $(x,0)$, e a soma de $\mathbb{N}$ é a primeira coordenada
daquele que se alcança. As duas operações que a iteração usa são a soma e o
produto de $\FF_2$, as únicas primitivas de um bit; o produto de $\mathbb{N}$
segue por deslocamento e soma.
\end{teorema}

\begin{proof}
Preservação: Lema~\ref{lem:carry}. Terminação, que é onde está a tese ---
com a medida escrita, porque é ela que faz o argumento:
\[
v(x):=\text{índice do bit não nulo mais baixo de }x\quad(x\neq0),
\qquad v(0):=+\infty .
\]
Seja $(a',b')=(a\oplus b,\,2(a\wedge b))$ o passo, com $b\neq0$. Como
$a\wedge b$ só tem bits onde \emph{ambos} têm, $v(a\wedge b)\ge v(b)$; e
duplicar acrescenta um: $v(b')=v(a\wedge b)+1\ge v(b)+1$. Logo $v$ do segundo
componente \textbf{cresce estritamente} a cada passo. Falta o tecto: se
$a,b<2^{n}$ então $a\oplus b<2^{n}$ e $2(a\wedge b)<2^{n+1}$, e a partir daí a
propriedade mantém-se --- nenhum componente passa de $2^{n+1}$. Um valor não
nulo abaixo de $2^{n+1}$ tem $v\le n$; como $v$ cresce de pelo menos $1$ por
passo a partir de $v(b)\ge0$, ao fim de no máximo $n+1$ passos já não há valor
não nulo possível, isto é, $b=0$.
\emph{Verifica:} \code{naturais.c} §NB7 --- nos $65\,536$ pares de $\Word_8$
o resultado bate com a definição recursiva de Peano
($n+0=n$, $n+S(m)=S(n+m)$, \code{naturais.h}), zero falhas, e o número máximo
de passos observado é $9=8+1$ --- a medida encontra exactamente o limite
$n+1$ da prova, e não um limite mais folgado.
$\square$
\end{proof}

\begin{corolario}[quanto vale o transporte]\label{cor:3n}
$a\oplus b=a+b$ exactamente quando $a\wedge b=0$. Em $\Word_8$ isso são
$3^{8}=6561$ dos $65\,536$ pares --- cada posição admite três estados
($0/0$, $1/0$, $0/1$) e o quarto é o que gera transporte.
\emph{A soma de $\FF_2$ é a soma de $\mathbb{N}$ com o transporte desligado.}
(\code{naturais.c} §NB6: $6561$ dos dois lados, e o $\iff$ verificado par a par.)
\end{corolario}

\medskip\noindent\textbf{A leitura.} A torre dá o \emph{conjunto} e a
\emph{reversibilidade}; o que ela não guarda é a memória entre posições. Em
$\FF_2$ cada bit opera isolado --- é isso que faz a operação bit a bit
(Thm.~\ref{thm:base}) --- e é isso, exactamente, que $\mathbb{N}$ acrescenta:
\[
\boxed{\;\mathbb{N}\big|_{\Word_8}\;=\;F_8\;+\;\text{transporte}\;}
\]
\medskip\noindent\textbf{E a ordem, em gaveta separada.} O
Thm.~\ref{thm:transporte} reconstrói a \emph{soma}, e só ela. A desigualdade
$a+b\ge a$ é uma propriedade da ordem usual de $\mathbb{N}$ --- que vive no
suporte enquanto conjunto de naturais, não nas operações do corpo --- e vale
para a soma reconstruída porque essa soma \emph{é} a soma usual, não porque o
teorema a produza. Do outro lado: $F_8$, com as suas próprias operações, não
admite ordem de corpo (§\ref{sec:escopo}). O mesmo conjunto de representantes
carrega as duas coisas, e elas não se demonstram uma à outra.

\colunas{\code{inteiros def:peano} (recebida do \code{inteiros.tex})}
        {Lema~\ref{lem:carry}; ponto fixo de $(\oplus,\wedge)$}
        {§NB6 ($6561=3^8$); §NB7 ($65\,536$ pares contra Peano, $\le9$ passos)}

%======================================================================
\section{A régua mínima é a do \emph{nim}}\label{sec:conway}

\noindent A nim-soma é o XOR e a nim-multiplicação é a operação que torna os
ordinais um corpo de característica $2$ (Conway, \emph{ONAG}, cap.~6). Dela
usam-se aqui duas regras, para $\Phi=2^{2^{k}}$ uma \emph{potência de Fermat}:
\[
\textbf{(N1)}\quad \Phi\otimes\Phi=\tfrac{3}{2}\Phi=\Phi\oplus\tfrac{\Phi}{2},
\qquad
\textbf{(N2)}\quad \Phi\otimes x=\Phi\,x\ \ \text{(produto ordinário)}\ \ \text{para }x<\Phi .
\]

\begin{teorema}[a torre é o nim]\label{thm:nim}
Admitidas \textbf{(N1)} e \textbf{(N2)}, o produto da torre da
Def.~\ref{def:dobra} com $\lambda_k=2^{2^{k}-1}=\Phi/2$ coincide com o
nim-produto em \textbf{todos} os andares.
\end{teorema}

\begin{proof}
Indução no andar. \emph{Base:} em $\{0,1\}$ o nim-produto é o AND, que é o
produto de $F_1$. \emph{Passo:} sejam $a=a_1\Phi\oplus a_0$ e
$b=b_1\Phi\oplus b_0$ com $a_i,b_i<\Phi$. Pela distributividade do nim,
\[
a\otimes b=(a_1\otimes b_1)\otimes(\Phi\otimes\Phi)
\;\oplus\;\bigl(a_1\otimes b_0\oplus a_0\otimes b_1\bigr)\otimes\Phi
\;\oplus\;a_0\otimes b_0 .
\]
Por \textbf{(N1)}, $(a_1\otimes b_1)\otimes(\Phi\oplus\Phi/2)
=(a_1\otimes b_1)\otimes\Phi\;\oplus\;(a_1\otimes b_1)\otimes(\Phi/2)$; e por
\textbf{(N2)} cada $\otimes\Phi$ é o deslocamento de $2^{k}$ bits. Separando a
parte alta da baixa:
\[
\text{alto}=a_1b_1\oplus a_1b_0\oplus a_0b_1,
\qquad
\text{baixo}=a_0b_0\oplus\lambda\,(a_1b_1),\qquad \lambda=\Phi/2,
\]
que é \emph{exactamente} a Def.~\ref{def:dobra}. Os produtos $a_ib_j$ são do
andar de baixo, onde as duas operações já coincidem por hipótese de indução.
$\square$
\end{proof}

\begin{observacao}[estatuto de cada peça, que é onde este resultado se podia inflar]
\textbf{(N1)} e \textbf{(N2)} são resultados de Conway --- \emph{citados}, não
demonstrados aqui. O que este documento acrescenta é a indução acima e, do
lado da medida, a verificação das duas regras \textbf{na própria torre}:
$\Phi\otimes\Phi=\Phi\oplus\Phi/2$ e $\Phi\otimes x=\Phi x$ para todo
$x<\Phi$, nos três andares, zero falhas (\code{naturais.c} §NB10). A ponte
para o \emph{nim} não é só citada: é medida do lado de cá.
\end{observacao}

\begin{proposicao}[verificação independente, e o seu escopo]\label{prop:mex}
O nim-produto definido pelo \emph{mex},
\[
a\otimes b=\operatorname{mex}\{\,a'\otimes b\oplus a\otimes b'\oplus a'\otimes b'
\;:\;a'<a,\;b'<b\,\},
\]
--- definição puramente combinatória, que não conhece $\sigma$, $\lambda$ nem
polinómio nenhum --- coincide com o produto da torre. \textbf{Escopo:} medido
exaustivamente em $F_4$ ($256$ pares, zero falhas) e verificado nos mesmos
pares dentro de $F_8$; o topo inteiro pelo mex \textbf{não} se varreu
($256^{4}$ candidatos), e isso não é o que sustenta o Thm.~\ref{thm:nim} --- o
que o sustenta é a indução, e esta proposição é o controlo por um caminho que
não partilha código com ela.
\end{proposicao}

\begin{proposicao}[o menor $\lambda$, medido]\label{prop:menor}
Nos três andares varridos, $\lambda_k=2^{2^{k}-1}$ é também o \textbf{menor}
$\lambda$ que serve (\code{naturais.c} §NB0). Para $k\ge3$ isto é
\emph{conjectura}: não está medido nem demonstrado. O Thm.~\ref{thm:nim} não
depende disto --- depende só de $\lambda_k=\Phi/2$.
\end{proposicao}

\medskip\noindent A leitura que isto sustenta, e só ela: os naturais que Conway
obteve como corpo pela combinatória de um jogo e os que a dobra dual constrói
pela álgebra são o mesmo objecto, e a razão comum é o \textbf{encaixe nos
prefixos} --- a propriedade que o Thm.~\ref{thm:encaixe} demonstra e que a
realização pelo polinómio do AES não tem (§\ref{sec:escopo}).

%======================================================================
\section{O escopo, e o gume}\label{sec:escopo}

\noindent\textbf{O que $F_8$ não é.} $F_8$ \textbf{não} é $\mathbb{N}$:
é finito; tem característica $2$ --- medido, $1+1=0$ em duas parcelas --- e
portanto não admite ordem de corpo (\code{algebrico}, §O byte). O que é
verdade, e é o que se mediu: o \emph{suporte} de $F_8$ é $\Word_8$ pela
identidade; os andares de baixo são os prefixos; e as duas operações de
$\mathbb{N}$ nesse suporte reconstroem-se das duas de $\FF_2$ mais o
transporte (Thm.~\ref{thm:transporte}). $\mathbb{N}$ inteiro não é o topo desta
torre --- a torre não pára em $8$, e o que pára em $8$ é a norma
(\code{algebrico}, Hurwitz), não o objecto.

\medskip\noindent\textbf{O gume: duas réguas para o mesmo conjunto.}
Comece-se por dizer o que \emph{não} é o resultado: quaisquer dois corpos de
$256$ elementos são isomorfos, e portanto nada se afirma sobre o corpo
\emph{abstracto}. O que se compara são \textbf{realizações} --- duas maneiras
de pôr as operações sobre o mesmo conjunto de representantes
$\{0,\ldots,255\}$.

O produto do AES (módulo $x^{8}+x^{4}+x^{3}+x+1$) é uma régua perfeitamente
válida: dá corpo, e todo não nulo tem inversa (medido: zero elementos sem
inversa). O que ela não tem é a filtração:
\[
2\otimes_{\mathrm{AES}}8=x\cdot x^{3}=x^{4}=16\notin\{0,\ldots,15\},
\qquad\text{enquanto}\qquad
2\otimes_{\mathrm{torre}}8=12\in\{0,\ldots,15\},
\]
e ao todo $176$ dos $256$ produtos de $\{0,\ldots,15\}$ saem do conjunto no
AES. (Um isomorfismo entre as duas realizações existe, claro --- o que ele não
faz é preservar os prefixos.) Logo:
\[
\boxed{\;\text{a REALIZAÇÃO pela dobra dual possui a filtração pelos prefixos}
\ \{0,\ldots,2^{2^{k}}-1\};\ \text{a do AES não.}\;}
\]
Não é o byte nem a característica $2$ que dão o encaixe --- é a torre. A régua
que se escolhe é parte do resultado (\code{naturais.c} §NB8), e este é
provavelmente o resultado estrutural do documento.

\medskip\noindent\textbf{O gume automático.} Três mutações no objecto,
compiladas e corridas: régua fixa em $1$ em vez de derivada
($4$ asserções caem); o alto do produto sem o cancelamento de Karatsuba
($3$ caem); o transporte sem deslocamento ($1$ cai). Nenhuma sobreviveu
silenciosa.

\medskip\noindent\textbf{Duas dívidas pagas.} \textbf{(i)} O teorema de
\textbf{Kummer} --- o número de transportes ao somar $a$ e $b$ em base $2$ é
$v_2\binom{a+b}{a}$ --- está agora medido: contar os transportes dá exactamente
$s_2(a)+s_2(b)-s_2(a+b)$, que por Legendre é essa valuação, nos $65\,536$ pares
de $\Word_8$ sem uma divergência (§NB12). O que o Cor.~\ref{cor:3n} contava
--- \emph{quantos} pares não têm transporte --- ganha o outro lado: \emph{quantos}
transportes tem cada par, e o número não é empírico, é a valuação $2$-ádica de um
binomial.
\textbf{(ii)} O \emph{menor} $\lambda$ ganhou mais um andar: no degrau que ergue
$F_8$ até $F_{16}$ (os $65\,536$ elementos), a busca devolve
$\lambda_3=128=\Phi/2$, e servem $128$ dos $256$ candidatos --- a mesma
repartição ao meio da Prop.~\ref{prop:regua}, pelo traço (§NB13). Continua
conjectura para $k\ge4$.

\medskip\noindent\textbf{Por medir}, e diz-se: as regras
\textbf{(N1)}/\textbf{(N2)} em geral --- aqui são citadas de Conway e medidas
só nos três andares (Thm.~\ref{thm:nim}); a verificação pelo mex tem o escopo
da Prop.~\ref{prop:mex}; e o teorema de Kummer --- o número de transportes ao somar $a$ e $b$ em base $2$ é
$v_2\binom{a+b}{a}$ --- que dá o Cor.~\ref{cor:3n} em forma fechada, e está medido em §NB12.

%======================================================================
\section{Referências}\label{sec:refs}

\subsection*{Na casa --- o que este paper recebe, e onde está}

\begin{center}\small
\begin{tabular}{@{}llp{7.2cm}@{}}
\toprule
onde & § & o que dá \\
\midrule
\code{papers/corpo\_algebrico.tex} & §O byte & $B=\sum b_k2^k$, posição $k$ à Lei $k$; a dobra $1\to2\to4\to8$; o produto módulo irredutível de grau oito; e o escopo de «não ordenado» \\
\code{papers/corpo\_algebrico.tex} & \code{thm:autossimilar} & o construtor $(K,m)\mapsto K[\sigma_m]$, $\sigma^2=m\sigma+1$: o mesmo código dá andar após andar; «o metal do andar 1 é a régua do andar 2» \\
\code{papers/corpo\_algebrico.tex} & \code{thm:gerador-primos} & a escada de Fermat $2^{2^{j}}$: $16\to17$, $256\to257$, $65536\to65537$ \\
\code{papers/inteiros.tex} & \code{sec:palavra8} & $\Word_8$, leitura cruzada exacta, wrap/saturação contados à parte \\
\code{papers/arquitetura.tex} & \code{sec:descida} & a cadeia $\Word_8\to E_{16}\to\FF_{127}\to\mathrm{GF}(2)$ \\
\code{tests/ortonormal.c} & §O0--§O5 & a base das oito leis é ortonormal (Gram $=I$); a coordenada é ler o bit; o preço (vectores isótropos) \\
\code{tests/umbit.c} & §B0--§B7 & $\FF_2$ nu; $\mathbb{P}^1(\FF_2)$; as oito leis como transformações do byte; o ciclo de período oito \\
\code{lib/corpo256.h} & --- & o byte como corpo pela régua do AES; $c6\_base(k)$ \\
\code{lib/gf2n.h} & --- & $\mathrm{GF}(2^n)=\FF_2[x]/(f_n)$; o Frobenius como dual \\
\code{lib/naturais.h} & --- & Peano P1--P5, soma e produto recursivos; $\Word_8$ \\
\code{lib/booleana.h} & --- & a ANF de Zhegalkin pela transformada de Möbius (involução) \\
\bottomrule
\end{tabular}
\end{center}

\noindent\textbf{Acrescentado por este paper:} \code{lib/binario.h} (a torre) e
\code{tests/naturais.c} (§NB0--§NB11).

\subsection*{Fora --- e o que cada uma responde}

\begin{itemize}\itemsep2pt
\item G.~Peano, \emph{Arithmetices principia, nova methodo exposita}, 1889;
H.~Grassmann, \emph{Lehrbuch der Arithmetik}, 1861 --- a definição recursiva
contra a qual §NB7 mede.
\item J.~H.~Conway, \emph{On Numbers and Games}, Academic Press, 1976,
cap.~6 --- o Corpo $\mathrm{On}_2$: os ordinais com nim-soma e nim-produto
formam um corpo algebricamente fechado de característica $2$, e os naturais
menores que $2^{2^{n}}$ formam subcorpo. \textbf{De lá vêm (N1) e (N2)}, as
duas regras das potências de Fermat de que a prova do Thm.~\ref{thm:nim}
depende --- e que este documento não demonstra, só mede na torre.
\item H.~W.~Lenstra Jr., \emph{Nim multiplication}, Séminaire de Théorie des
Nombres de Bordeaux, 1977--78 --- a estrutura do mesmo corpo.
\item J.~H.~Conway, N.~J.~A.~Sloane, \emph{Lexicographic codes:
error-correcting codes from game theory}, IEEE Trans.\ Inform.\ Theory 32
(1986) --- a mesma aritmética a construir códigos, que é a ponte para
$E_{16}$ da \code{arquitetura sec:descida}.
\item E.~E.~Kummer, J.~reine angew.\ Math.\ 44 (1852) --- o número de
transportes na soma em base $p$ é $v_p\binom{a+b}{a}$: a forma fechada do
Cor.~\ref{cor:3n}, \emph{por medir}.
\item A.~Satoh \emph{et al.}, ASIACRYPT 2001; D.~Canright, \emph{A very
compact S-box for AES}, CHES 2005 --- a mesma torre
$F_1\to F_2\to F_4\to F_8$ (na notação por cardinais daqueles autores,
$\FF_2\to\FF_4\to\FF_{16}\to\FF_{256}$) usada para reduzir hardware: a
construção é conhecida na engenharia; o que aqui se acrescenta é o
\emph{encaixe nos prefixos} como propriedade e o transporte como o que falta
para $\mathbb{N}$.
\item J.~Daemen, V.~Rijmen, \emph{The Design of Rijndael}, Springer 2002 ---
a régua $x^8+x^4+x^3+x+1$ do gume de §\ref{sec:escopo}.
\item R.~D.~Schafer, \emph{An Introduction to Nonassociative Algebras}, 1966;
J.~Baez, \emph{The Octonions}, Bull.\ AMS 39 (2002) --- Cayley--Dickson,
$A_{n+1}=A_n\oplus A_n$ com conjugação: o mesmo molde $T+T^{*}$ em
característica $0$, onde a norma pára em $8$.
\item I.~Zhegalkin, 1927; M.~H.~Stone, 1936 --- a lógica é $\FF_2$
(\code{booleana.h}).
\item D.~E.~Knuth, \emph{TAOCP} vol.~2 §4.3.1 --- a adição em base $2$ e o
transporte.
\end{itemize}

%======================================================================
\section*{Síntese}

\[
\boxed{
\begin{array}{c}
F_1\;\xrightarrow{\;+\sigma_0\;}\;F_2\;\xrightarrow{\;+\sigma_1\;}\;F_4
\;\xrightarrow{\;+\sigma_2\;}\;F_8\;=\;\Word_8,
\qquad F_w=\mathrm{GF}(2^{w})\\[3pt]
\sigma_j^{2}=\sigma_j+\lambda_j,\qquad \lambda_j=2^{2^{j}-1},\qquad
e_k=\textstyle\prod_{k_j=1}\sigma_j=2^{k}\\[3pt]
\mathbb{N}\big|_{\Word_8}=F_8+\text{transporte},\qquad
a+b=(a\oplus b)+2(a\wedge b)
\end{array}}
\]

\begin{itemize}\itemsep2pt
\item[\textbf{0}] $x\oplus x=0$ imediato; a subtracção é total no chão; só o $0$ fica sem inversa.
\item[\textbf{1}] degenera ($-x=x$); o dual muda de sítio: é a segunda raiz, $x^{*}$.
\item[\textbf{2}] $x^{**}=x$; $x^{*}$ é automorfismo dos dois lados e fixa o andar de baixo.
\item[\textbf{3}] degenera em $\{-1,+1\}$: sem cúspide, a torre dobra sem excepção.
\item[\textbf{4}] $F_{2w}=F_w\oplus\sigma F_w$ --- o molde do \code{inteiros.tex} sem quociente.
\item[\textbf{5}] degenera: $x^2=1$ tem uma só raiz; não há $i$ neste corpo.
\item[\textbf{6}] a base dos produtos dos geradores é ortonormal, e por isso a coordenada é ler o bit.
\item[\textbf{7}] encaixe por prefixos: os andares ligam-se sem se fundirem.
\end{itemize}

\medido{\code{tests/naturais.c} §NB0--§NB13 ($14{:}0$): torre e inversas §NB0;
soma involutiva §NB1; dual §NB2; encaixe §NB3 ($276$ pares);
Frobenius §NB4; trial §NB5 ($11$/$0$/$11$); transporte §NB6 ($3^{8}=6561$);
Peano §NB7 ($65\,536$ pares, $\le9$ passos); gume das duas réguas §NB8
($176$ produtos saem no AES); escopo §NB9; \emph{nim} §NB10 --- mex em $F_4$,
e as regras (N1) $\Phi\otimes\Phi=\Phi\oplus\Phi/2$ e (N2)
$\Phi\otimes x=\Phi x$ medidas na torre, zero falhas; base
ortonormal §NB11; \textbf{Kummer} §NB12 ($65\,536$ pares, zero divergências,
máximo $8$ transportes); $\lambda_3=128$ §NB13. Gume: três mutações, $4/3/1$ asserções caídas.
Recebido: \code{tests/ortonormal.c} §O0; \code{tests/umbit.c} §B0--§B7.}

\medskip\noindent
\code{papers/inteiros.tex} · \code{papers/corpo\_algebrico.tex} ·
\code{papers/arquitetura.tex} · \code{lib/binario.h}.

\vfill
{\footnotesize\color{cinza}\noindent Proprietário --- uso mediante contrato e pagamento (ver \texttt{LICENSE}).\par}

\end{document}
